% ============================================================
% 基礎演習Ⅱ　産学連携プロジェクト　グループ活動記録シート
%   第6〜12回のグループディスカッションごとに1枚ずつ記入し、
%   第13回（プロジェクトの総括）で全体を振り返る。
%   デザインは各回ワークシート（../worksheets/preamble.tex）と共通。
%   ビルド：make（Makefile を参照）
% ============================================================
% 共通のプリアンブル（worksheets.tex と session.tex から読み込む）
\documentclass[paper=a4paper,fontsize=10pt,line_length=48zw,number_of_lines=44,baselineskip=1.6zw]{jlreq}

\usepackage{xcolor}
\usepackage[most]{tcolorbox}
\usepackage{tabularray}
\usepackage{enumitem}
\usepackage{needspace}
\usepackage{fancyhdr}
\usepackage[hidelinks]{hyperref}

% ------------------------------------------------------------
% 色（サイトの assets/css/style.css に合わせる）
% ------------------------------------------------------------
\definecolor{ink}{HTML}{1D2330}
\definecolor{inktwo}{HTML}{4A5261}
\definecolor{inkthree}{HTML}{757D8A}
\definecolor{rulegray}{HTML}{B9BEB4}
\definecolor{papertwo}{HTML}{ECEEE8}
\definecolor{pone}{HTML}{2747C9}  \definecolor{ponetint}{HTML}{DCE2F9}
\definecolor{ptwo}{HTML}{D68A17}  \definecolor{ptwotint}{HTML}{F6E4C3}
\definecolor{pthree}{HTML}{6E7684}\definecolor{pthreetint}{HTML}{E0E3E4}
\definecolor{bad}{HTML}{B3261E}   \definecolor{badtint}{HTML}{F9DEDA}
\colorlet{phase}{pone}
\colorlet{phasetint}{ponetint}

% ------------------------------------------------------------
% ページスタイル
% ------------------------------------------------------------
\newcommand*\sessionmark{}
\pagestyle{fancy}
\fancyhf{}
\renewcommand{\headrulewidth}{0pt}
\fancyfoot[L]{\sffamily\scriptsize\color{inkthree}基礎演習Ⅱ　2026年度後期　ワークシート}
\fancyfoot[R]{\sffamily\scriptsize\color{inkthree}\sessionmark\quad\thepage}
\fancypagestyle{cover}{\fancyhf{}}

\setlength{\parindent}{0pt}
\setlist{nosep,leftmargin=1.5em}

% ------------------------------------------------------------
% 氏名欄（記入例・模範解答版では answers-preamble.tex で中身を入れる）
\newcommand*\wsid{}
\newcommand*\wsname{}
\newcommand*\wsgroup{}
\newcommand\fieldline[2]{\makebox[#1][l]{\rlap{\hspace{1mm}#2}\hrulefill}}

% 各回の見出し
%   \session{回}{フェーズ 1|2|3}{日付}{題目}{タグ}{グループ欄 yes|no}
% ------------------------------------------------------------
\newcommand\session[6]{%
  \clearpage
  \ifcase#2\or
    \colorlet{phase}{pone}\colorlet{phasetint}{ponetint}\def\phasename{探究の基礎}\or
    \colorlet{phase}{ptwo}\colorlet{phasetint}{ptwotint}\def\phasename{実践プロジェクト}\or
    \colorlet{phase}{pthree}\colorlet{phasetint}{pthreetint}\def\phasename{キャリア講習}\fi
  \renewcommand*\sessionmark{第#1回　#4}%
  \phantomsection\addcontentsline{toc}{section}{\texorpdfstring{第#1回　#4\hfill{\small #3}}{第#1回　#4}}%
  \begin{tcolorbox}[enhanced,colback=phase,colframe=phase,arc=2mm,boxrule=0pt,
      left=4mm,right=4mm,top=2.5mm,bottom=2.5mm,fontupper=\color{white}\sffamily]
    {\small 第#1回　#3　\textbar　\phasename\if\relax\detokenize{#5}\relax\else　\textbar　#5\fi}\par
    \vspace{0.6mm}{\LARGE\bfseries #4}
  \end{tcolorbox}
  \vspace{-1mm}
  {\sffamily\small\color{inktwo}
  学籍番号\fieldline{30mm}{\wsid}\quad 氏名\fieldline{42mm}{\wsname}%
  \ifx\relax#6\relax\else\quad グループ\fieldline{28mm}{\wsgroup}\fi}\par
  \vspace{2mm}%
}

% ねらい
\newcommand\goals[1]{%
  \begin{tcolorbox}[enhanced,colback=phasetint,colframe=phasetint,arc=1.5mm,boxrule=0pt,
      left=3mm,right=3mm,top=1.5mm,bottom=1.5mm,fontupper=\small]
    {\sffamily\bfseries\color{phase}この回のねらい}\par
    \begin{itemize}[label=\textcolor{phase}{・},leftmargin=1.2em]#1\end{itemize}
  \end{tcolorbox}}

% ワークの見出し　\work[必要な高さ]{ラベル}{題}
\newcommand\work[3][36mm]{%
  \par\needspace{#1}\vspace{2.5mm}%
  {\sffamily\bfseries\color{phase}■ #2}\enspace{\sffamily\bfseries #3}\par\vspace{1mm}}

% 指示文
\newcommand\instr[1]{{\small\color{inktwo}#1}\par\vspace{1mm}}

% 記入欄　\answerbox[欄内の小見出し]{高さ}
\newcommand\answerbox[2][]{%
  \begin{tcolorbox}[enhanced,colback=white,colframe=rulegray,boxrule=0.5pt,arc=1.2mm,
      height=#2,left=2mm,right=2mm,top=1.2mm,bottom=1mm,
      fontupper=\sffamily\scriptsize\color{inkthree},before skip=1mm,after skip=1.5mm]#1\end{tcolorbox}}

% 横並びの記入欄　\twoboxes{見出し1}{見出し2}{高さ}
\newcommand\twoboxes[3]{%
  \par\noindent
  \begin{minipage}[t]{0.49\linewidth}\answerbox[#1]{#3}\end{minipage}\hfill
  \begin{minipage}[t]{0.49\linewidth}\answerbox[#2]{#3}\end{minipage}\par}

% 方眼の記入欄（スケッチ用）
\newcommand\gridbox[2][]{%
  \begin{tcolorbox}[enhanced,colback=white,colframe=rulegray,boxrule=0.5pt,arc=1.2mm,
      height=#2,left=2mm,top=1.2mm,fontupper=\sffamily\scriptsize\color{inkthree},
      underlay={\begin{tcbclipinterior}\draw[step=5mm,black!12,very thin]
        (interior.south west) grid (interior.north east);\end{tcbclipinterior}}]#1\end{tcolorbox}}

% 罫線　\writelines{行数}
\newcommand\writelines[1]{%
  \par\foreach\n in {1,...,#1}{\nointerlineskip\noindent
    {\color{rulegray}\vbox to 7.5mm{\vfil\hrule width\linewidth}}\par}\vspace{1.5mm}}

% 1行の記入欄　\fillin{見出し}
\newcommand\fillin[1]{%
  \par\noindent{\sffamily\small #1}\enspace{\color{rulegray}\leaders\hrule\hfill}\par\vspace{2.5mm}}

% プロンプト例
\newtcolorbox{promptbox}[1][]{enhanced,colback=papertwo,colframe=papertwo,boxrule=0pt,
  arc=1mm,leftrule=1.2mm,colbacktitle=papertwo,
  borderline west={1.2mm}{0pt}{phase},left=3mm,right=2mm,top=1mm,bottom=1mm,
  fontupper=\sffamily\footnotesize\color{inktwo},before skip=1mm,after skip=2mm,
  title={\sffamily\scriptsize\bfseries\color{phase}プロンプト例　#1},
  attach title to upper={\par\vspace{0.5mm}}}

% 注意
\newtcolorbox{caution}[1]{enhanced,colback=badtint,colframe=badtint,boxrule=0pt,arc=1mm,
  left=3mm,right=3mm,top=1mm,bottom=1mm,fontupper=\footnotesize,
  before upper={{\sffamily\bfseries\color{bad}#1}\par}}

% チェック項目
\newenvironment{checks}{\begin{itemize}[label={□},leftmargin=1.5em,itemsep=0.8mm]\small}{\end{itemize}}

% 次回までに
\newenvironment{nexttime}{%
  \par\needspace{24mm}\vspace{2.5mm}%
  \begin{tcolorbox}[enhanced,colback=white,colframe=phase,boxrule=0.6pt,arc=1.5mm,
      left=3mm,right=3mm,top=1.5mm,bottom=1.5mm,fontupper=\small]
    {\sffamily\bfseries\color{phase}次回までに}\par\vspace{0.5mm}
    \begin{itemize}[label={□},leftmargin=1.5em,itemsep=0.6mm]}%
  {\end{itemize}\end{tcolorbox}}

% AI活用記録（その回の分）
\newcommand\aimemo{%
  \par\needspace{50mm}\vspace{2.5mm}%
  {\sffamily\bfseries\color{phase}■ AI活用記録（この回の分）}\par\vspace{1mm}
  \begin{wstab}{colspec={Q[l,wd=30mm,m] X[l,m]},row{1-Z}={ht=9mm}}
    使用した場面 & \\
    主なプロンプト & \\
    出力をどう扱ったか\newline{\scriptsize 採用・修正・不採用と理由} & \\
    検証したこと\newline{\scriptsize 何で確かめたか} & \\
  \end{wstab}\par}

% 表の共通スタイル
\NewTblrEnviron{wstab}
\SetTblrInner[wstab]{hlines={0.4pt,rulegray},vlines={0.4pt,rulegray},
  cells={font=\small},rowsep=2pt,colsep=4pt}
\NewTblrEnviron{wstabh}
\SetTblrInner[wstabh]{hlines={0.4pt,rulegray},vlines={0.4pt,rulegray},
  cells={font=\small},rowsep=2pt,colsep=4pt,
  row{1}={bg=phasetint,font=\sffamily\footnotesize\bfseries,ht=0pt}}

% ペルソナカード（第6回）
\newcommand\personacard[1]{%
\begin{wstab}{colspec={Q[l,wd=26mm,m] X[l,m]},rows={ht=7.5mm}}
  \SetCell[c=2]{bg=phasetint,font=\sffamily\footnotesize\bfseries}ペルソナ#1　名前： & \\
  年齢・職業 & \\ 家族構成 & \\ 居住地の特徴 & \\ 来店の目的 & \\ 重視する点 & \\
  不安や面倒 & \\ 休日の過ごし方 & \\
  \SetRow{ht=14mm}本人の言葉\newline\scriptsize 店舗に求めること & \\
\end{wstab}}

% 見出しセル（左列を色付きにする用）
\newcommand\hc[1]{\SetCell{bg=phasetint,font=\sffamily\footnotesize\bfseries}#1}


\colorlet{phase}{ptwo}
\colorlet{phasetint}{ptwotint}
\fancyfoot[L]{\sffamily\scriptsize\color{inkthree}基礎演習Ⅱ　2026年度後期　グループ活動記録シート}

% ------------------------------------------------------------
% ページの見出し　\loghead{上の行}{題目}
% ------------------------------------------------------------
\newcommand\loghead[2]{%
  \clearpage
  \renewcommand*\sessionmark{#2}%
  \begin{tcolorbox}[enhanced,colback=phase,colframe=phase,arc=2mm,boxrule=0pt,
      left=4mm,right=4mm,top=2.5mm,bottom=2.5mm,fontupper=\color{white}\sffamily]
    {\small #1}\par
    \vspace{0.6mm}{\LARGE\bfseries #2}
  \end{tcolorbox}
  \vspace{-1mm}
  {\sffamily\small\color{inktwo}
  学籍番号\fieldline{30mm}{}\quad 氏名\fieldline{42mm}{}\quad グループ\fieldline{28mm}{}}\par
  \vspace{1mm}}

% 小見出し（\work より詰める）
\newcommand\logsec[2]{%
  \par\vspace{1.2mm}%
  {\sffamily\bfseries\color{phase}■ #1}\enspace{\sffamily\bfseries #2}\par\vspace{0.8mm}}

% 選択肢の□を並べる
\newcommand\opt[1]{□\,#1\quad}

% 毎回の記録　\logpage{回}{日付}{題目}
\newcommand\logpage[3]{%
  \loghead{グループ活動記録　\textbar　第#1回　#2}{#3}

  \logsec{1}{出席}
  {\small 自分：\enspace\opt{出席}\opt{遅刻}\opt{早退}\opt{欠席}
    {\color{inktwo}\footnotesize → 出席以外のときは、次の2行も書く}}\par
  \vspace{1.2mm}
  \fillin{事前の連絡\enspace □\,した\enspace □\,していない\quad 共有してもらった内容と相手}
  \fillin{追いつくためにしたこと／グループのためにできること}
  {\small メンバーの出欠}\enspace{\footnotesize\color{inktwo}番号は1ページの名簿の番号。○ 出席　△ 遅刻・早退　× 欠席}\par
  \vspace{0.8mm}
  \begin{wstab}{colspec={Q[l,wd=16mm,m] X[c,m] X[c,m] X[c,m] X[c,m] X[c,m] X[c,m]},
      row{1}={bg=phasetint,font=\sffamily\footnotesize\bfseries},row{2}={ht=7mm}}
    番号 & ① & ② & ③ & ④ & ⑤ & ⑥ \\
    出欠 & & & & & & \\
  \end{wstab}\par\vspace{1.5mm}
  \fillin{欠席者への共有（誰が・何で・何を）}

  \logsec{2}{ディスカッションの成果}
  \fillin{今日の目的・議題}
  \answerbox[グループで決まったこと・できあがったもの（成果物）]{19mm}
  \twoboxes{まだ決まっていないこと・次回に持ち越すこと}{意見が分かれた点と、どう決めたか}{14mm}

  \logsec{3}{自分の役割}
  {\small 今日の役割：\enspace\opt{司会・進行}\opt{記録}\opt{調べ役}\opt{AI操作}\opt{資料作成}\opt{発表}\opt{時間管理}\opt{その他（\hspace{10mm}）}}\par
  \vspace{0.5mm}
  \answerbox[実際にやったこと・グループへの貢献（具体的に）]{14mm}
  \par\noindent{\sffamily\small 役割を果たせた度合い\enspace 1\enspace・\enspace 2\enspace・\enspace 3\enspace・\enspace 4\enspace・\enspace 5\quad 理由}\enspace{\color{rulegray}\leaders\hrule\hfill}\par\vspace{1.5mm}

  \logsec{4}{気づき}
  \twoboxes{グループの進め方でうまくいったこと}{次に改善したいこと}{14mm}
  \answerbox[自分自身について気づいたこと（得意・苦手、メンバーから学んだこと）]{12mm}

  \logsec{5}{次回までの分担}
  \begin{wstabh}{colspec={X[l,m] Q[l,wd=30mm,m] Q[l,wd=22mm,m] Q[c,wd=8mm,m]},row{2-Z}={ht=6mm}}
    やること & 担当 & 期限 & 済 \\
    & & & □ \\ & & & □ \\ & & & □ \\
  \end{wstabh}\par
}

\begin{document}

% ============================================================
% 1ページ目：使い方・名簿・ルール
% ============================================================
\renewcommand*\sessionmark{使い方・名簿・ルール}
\begin{tcolorbox}[enhanced,colback=phase,colframe=phase,arc=2mm,boxrule=0pt,
    left=4mm,right=4mm,top=2.5mm,bottom=2.5mm,fontupper=\color{white}\sffamily]
  {\small 基礎演習Ⅱ　2026年度後期　\textbar　実践プロジェクト（第6〜13回）}\par
  \vspace{0.6mm}{\LARGE\bfseries 産学連携プロジェクト　グループ活動記録シート}
\end{tcolorbox}
\vspace{-1mm}
{\sffamily\small\color{inktwo}
学籍番号\fieldline{30mm}{}\quad 氏名\fieldline{42mm}{}\quad グループ\fieldline{28mm}{}}\par
\vspace{2mm}

\begin{tcolorbox}[enhanced,colback=phasetint,colframe=phasetint,arc=1.5mm,boxrule=0pt,
    left=3mm,right=3mm,top=1.5mm,bottom=1.5mm,fontupper=\small]
  {\sffamily\bfseries\color{phase}このシートの使い方}\par
  \begin{itemize}[label=\textcolor{phase}{・},leftmargin=1.2em]
    \item グループで話し合った回ごとに、授業の最後の5分で1枚（3〜9ページ）を\textbf{個人で}記入する。授業外に集まったときなどは予備のページ（10・11ページ）を使う。
    \item 毎回、①その日の出欠、②ディスカッションで得られた成果、③自分の役割と実際にやったこと、④気づき、を書く。成果は「話した」ではなく「何が決まったか、何ができたか」で書く。
    \item 欠席・遅刻・早退した回も記入する。メンバーから共有を受けた内容と、追いつくためにしたことを書く。
    \item 記入したら、2ページの「出席と役割の記録」に1行で転記する。第13回（プロジェクトの総括）で最後のページを書き、全体を振り返る。
    \item 他のメンバーの出欠は、責めるためではなく、欠席者への共有と分担の見直しに使う。
  \end{itemize}
\end{tcolorbox}

\work[50mm]{名簿}{グループのメンバー（以後のページでは番号で書く）}
\begin{wstabh}{colspec={Q[c,wd=10mm,m] X[2,l,m] X[2,l,m] X[3,l,m]},row{2-Z}={ht=8mm}}
  番号 & 氏名 & 最初の役割 & 連絡方法・メモ \\
  ① & & & \\ ② & & & \\ ③ & & & \\ ④ & & & \\ ⑤ & & & \\ ⑥ & & & \\
\end{wstabh}

\work[50mm]{ルール}{グループで決めた約束}
\begin{wstab}{colspec={Q[l,wd=46mm,m] X[l,m]},row{1-Z}={ht=14mm}}
  \hc{欠席・遅刻の連絡\newline 誰に・いつまでに・何で} & \\
  \hc{欠席者への共有のしかた\newline 誰が・何で・いつまでに} & \\
  \hc{授業外の作業の進め方\newline 締め切り・集まり方} & \\
  \hc{意見が分かれたときの決め方} & \\
\end{wstab}

% ============================================================
% 2ページ目：出席と役割の記録の一覧
% ============================================================
\loghead{グループ活動記録　\textbar　第6〜12回の一覧}{出席と役割の記録}

\instr{毎回の記録を書いたら、ここに1行で転記する。出欠は ○ 出席　△ 遅刻・早退　× 欠席。欠席したメンバーは名簿の番号で書く。第13回の振り返りは、この一覧をもとに書く。}
\begin{wstabh}{colspec={Q[c,wd=9mm,m] Q[c,wd=13mm,m] Q[c,wd=11mm,m] Q[c,wd=15mm,m] X[2,l,m] X[3,l,m] X[3,l,m]},row{2-Z}={ht=19mm}}
  回 & 日付 & 自分の\newline 出欠 & 欠席した\newline メンバー & 自分の役割 & グループの成果（一言） & 気づき（一言） \\
  6 & 11/4 & & & & & \\
  7 & 11/11 & & & & & \\
  8 & 11/18 & & & & & \\
  9 & 11/25 & & & & & \\
  10 & 12/2 & & & & & \\
  11 & 12/9 & & & & & \\
  12 & 1/13 & & & & & \\
  予備 & & & & & & \\
  予備 & & & & & & \\
\end{wstabh}

% ============================================================
% 3〜9ページ：毎回の記録
% ============================================================
\logpage{6}{11月4日（水）}{産学連携キックオフ}
\logpage{7}{11月11日（水）}{コンセプト立案}
\logpage{8}{11月18日（水）}{店舗設計・空間構成}
\logpage{9}{11月25日（水）}{動線とUXのデザイン}
\logpage{10}{12月2日（水）}{提案のブラッシュアップ}
\logpage{11}{12月9日（水）}{プレゼン資料作成}
\logpage{12}{2027年1月13日（水）}{最終プレゼン発表会}

% 10・11ページ：予備（授業外の集まりや、日程が変わった回の用）
\logpage{\hspace{6mm}}{\hspace{8mm}月\hspace{8mm}日（\hspace{5mm}）\quad □ 授業内\enspace □ 授業外}{予備（追加の話し合い用）}
\logpage{\hspace{6mm}}{\hspace{8mm}月\hspace{8mm}日（\hspace{5mm}）\quad □ 授業内\enspace □ 授業外}{予備（追加の話し合い用）}

% ============================================================
% 12ページ：プロジェクト全体の振り返り（第13回）
% ============================================================
\loghead{グループ活動記録　\textbar　第13回　2027年1月20日（水）}{プロジェクト全体の振り返り}

\logsec{1}{出席の集計}
{\small 第6〜12回（全7回）のうち\quad 出席\fieldline{12mm}{}回\quad 遅刻・早退\fieldline{12mm}{}回\quad 欠席\fieldline{12mm}{}回}\par
\vspace{1mm}
{\footnotesize\color{inktwo}2ページの「出席と役割の記録」から数える。授業外の集まりへの参加は\fieldline{12mm}{}回}\par

\logsec{2}{自分の役割を振り返る}
\begin{wstab}{colspec={Q[l,wd=50mm,m] X[l,m]},row{1-Z}={ht=14mm}}
  \hc{自分が担った役割は、回を追ってどう変わったか} & \\
  \hc{グループの成果のうち、自分が最も貢献した部分と、その具体的な行動} & \\
  \hc{果たせなかった役割、もっとできたこと。その理由} & \\
  \hc{欠席や作業の遅れが出たとき、グループはどう対応したか。自分は何をしたか} & \\
\end{wstab}

\logsec{3}{気づきをまとめる}
\begin{wstab}{colspec={Q[l,wd=50mm,m] X[l,m]},row{1-Z}={ht=14mm}}
  \hc{グループで成果を出すために大事だと気づいたこと} & \\
  \hc{自分の強みと課題（チームの中での自分）} & \\
  \hc{次にグループで活動するとき、自分が意識すること} & \\
\end{wstab}

\logsec{4}{メンバーの良かった行動}
{\footnotesize\color{inktwo}名簿の番号ごとに、助けられたこと・見習いたいことを具体的に書く（自分の欄は空けてよい）}\par\vspace{0.8mm}
\begin{wstab}{colspec={Q[c,wd=10mm,m] X[l,m]},row{1-Z}={ht=6.5mm}}
  ① & \\ ② & \\ ③ & \\ ④ & \\ ⑤ & \\ ⑥ & \\
\end{wstab}

\end{document}
